% GENERATED FILE. Edit canonical lessons, then run npm run generate:books.
% canonical-source-hash: fnv1a64:a3055099770ae7b5
% canonical-lessons: BN-C342-ukil, BN-C342-bichar-kora, BN-C342-ray, BN-C342-nalish, BN-C342-mimangsa-kora

\chapter{In Court}
\label{ch:bn-in-court}
\hlchaptermodality{342}

\begin{chapteropening}
\textbf{By the end of this chapter:} \emph{I can say who heard a case and what was decided.}

Complete the last lesson of chapter 342: I can say who heard a case and what was decided.
\end{chapteropening}

\section[\texorpdfstring{\bn{উকিল} (ukil)}{উকিল (ukil)}]{\bn{উকিল} (ukil) — a lawyer}

\label{lesson:BN-C342-ukil}



\begin{quote}
Before the new one: say the Bengali for corruption, then the Bengali for an injustice.
\end{quote}

\subsection*{You'll want to know: \bn{উকিল}}
\textbf{\bn{উকিল}} — \emph{ukil} — \textquotedblleft{}a lawyer\textquotedblright{}.

\begin{quote}
\bn{আমাদের} \bn{একজন} \bn{ভালো} \bn{উকিল} \bn{দরকার।}
\end{quote}

\begin{quote}
\emph{āmāder ekjôn bhālo ukil dôrkār.}
\end{quote}

\begin{quote}
We need a good lawyer.
\end{quote}

\begin{cousinweb}
\textbf{\bn{উকিল}} (\emph{ukil}) is usually traced to Arabic \emph{wakīl}, one who acts for another, through Persian. The lawyer is the one who speaks for you.
\end{cousinweb}

\begin{grammarlens}[title={What you've built}]
One more word for giving reasons and opinions.
\end{grammarlens}

\subsection*{Your turn}
\begin{itemize}
\raggedright
  \item \emph{Say it:} \emph{ukil}
  \item \emph{Say it:} \emph{ukil}, once more
  \item \emph{Say it:} \emph{āmāder ekjôn bhālo ukil dôrkār.} — We need a good lawyer.
  \item \emph{Recall:} say the Bengali for corruption, then the Bengali for an injustice, then say \emph{ukil} again
\end{itemize}

\begin{tcolorbox}[breakable,colback=teal!4,colframe=teal!35!black,title={Before you move on}]
What does \bn{উকিল} mean? (\textquotedblleft{}A lawyer\textquotedblright{}.) Say it once more.
\end{tcolorbox}
\section[\texorpdfstring{\bn{বিচার} \bn{করা} (bichār kôrā)}{বিচার করা (bichār kôrā)}]{\bn{বিচার} \bn{করা} (bichār kôrā) — to try a case; to judge}

\label{lesson:BN-C342-bichar-kora}



\begin{quote}
Before the new one: say the Bengali for an injustice, then the Bengali for a lawyer.
\end{quote}

\subsection*{You'll want to know: \bn{বিচার} \bn{করা}}
\textbf{\bn{বিচার} \bn{করা}} — \emph{bichār kôrā} — \textquotedblleft{}to try a case; to judge\textquotedblright{}.

\begin{quote}
\bn{আদালত} \bn{এই} \bn{মামলার} \bn{বিচার} \bn{করবে।}
\end{quote}

\begin{quote}
\emph{ādālôt ei māmlār bichār kôrbe.}
\end{quote}

\begin{quote}
The court will try this case.
\end{quote}

Past (he/she …): \bn{তিনি} \bn{মামলাটার} \bn{বিচার} \bn{করলেন}, \emph{tini māmlāṭār bichār kôrlen}.

\begin{cousinweb}
\textbf{\bn{বিচার} \bn{করা}} (\emph{bichār kôrā}) is Sanskrit \emph{vicāra}, consideration or examination (Monier-Williams), with \bn{করা}. The judge, \bn{বিচারক}, is the one who does it.
\end{cousinweb}

\begin{grammarlens}[title={What you've built}]
One more word for giving reasons and opinions.
\end{grammarlens}

\subsection*{Your turn}
\begin{itemize}
\raggedright
  \item \emph{Say it:} \emph{bichār kôrā}
  \item \emph{Say it:} \emph{bichār kôrā}, once more
  \item \emph{Say it:} \emph{ādālôt ei māmlār bichār kôrbe.} — The court will try this case.
  \item \emph{Recall:} say the Bengali for an injustice, then the Bengali for a lawyer, then say \emph{bichār kôrā} again
\end{itemize}

\begin{tcolorbox}[breakable,colback=teal!4,colframe=teal!35!black,title={Before you move on}]
What does \bn{বিচার} \bn{করা} mean? (\textquotedblleft{}To try a case; to judge\textquotedblright{}.) Say it once more.
\end{tcolorbox}
\section[\texorpdfstring{\bn{রায়} (rāy)}{রায় (rāy)}]{\bn{রায়} (rāy) — a verdict; a judgment}

\label{lesson:BN-C342-ray}



\begin{quote}
Before the new one: say the Bengali for a lawyer, then the Bengali for to try a case.
\end{quote}

\subsection*{You'll want to know: \bn{রায়}}
\textbf{\bn{রায়}} — \emph{rāy} — \textquotedblleft{}a verdict; a judgment\textquotedblright{}.

\begin{quote}
\bn{আদালত} \bn{কাল} \bn{রায়} \bn{দেবে।}
\end{quote}

\begin{quote}
\emph{ādālôt kāl rāy debe.}
\end{quote}

\begin{quote}
The court will give its verdict tomorrow.
\end{quote}

\begin{cousinweb}
\textbf{\bn{রায়}} (\emph{rāy}) is Persian \emph{rāy}, an opinion or judgment, usually traced to Arabic \emph{raʾy}. A court gives it with \bn{দেওয়া}.
\end{cousinweb}

\begin{grammarlens}[title={What you've built}]
One more word for giving reasons and opinions.
\end{grammarlens}

\subsection*{Your turn}
\begin{itemize}
\raggedright
  \item \emph{Say it:} \emph{rāy}
  \item \emph{Say it:} \emph{rāy}, once more
  \item \emph{Say it:} \emph{ādālôt kāl rāy debe.} — The court will give its verdict tomorrow.
  \item \emph{Recall:} say the Bengali for a lawyer, then the Bengali for to try a case, then say \emph{rāy} again
\end{itemize}

\begin{tcolorbox}[breakable,colback=teal!4,colframe=teal!35!black,title={Before you move on}]
What does \bn{রায়} mean? (\textquotedblleft{}A verdict; a judgment\textquotedblright{}.) Say it once more.
\end{tcolorbox}
\section[\texorpdfstring{\bn{নালিশ} (nālish)}{নালিশ (nālish)}]{\bn{নালিশ} (nālish) — a formal complaint}

\label{lesson:BN-C342-nalish}



\begin{quote}
Before the new one: say the Bengali for to try a case, then the Bengali for a verdict.
\end{quote}

\subsection*{You'll want to know: \bn{নালিশ}}
\textbf{\bn{নালিশ}} — \emph{nālish} — \textquotedblleft{}a formal complaint\textquotedblright{}.

\begin{quote}
\bn{আমি} \bn{থানায়} \bn{নালিশ} \bn{জানালাম।}
\end{quote}

\begin{quote}
\emph{āmi thānāy nālish jānālām.}
\end{quote}

\begin{quote}
I lodged a complaint at the police station.
\end{quote}

\begin{cousinweb}
\textbf{\bn{নালিশ}} (\emph{nālish}) is usually traced to Persian \emph{nālish}, a lament or complaint. It is made to someone in authority, with \bn{জানানো} or \bn{করা}.
\end{cousinweb}

\begin{grammarlens}[title={What you've built}]
One more word for giving reasons and opinions.
\end{grammarlens}

\subsection*{Your turn}
\begin{itemize}
\raggedright
  \item \emph{Say it:} \emph{nālish}
  \item \emph{Say it:} \emph{nālish}, once more
  \item \emph{Say it:} \emph{āmi thānāy nālish jānālām.} — I lodged a complaint at the police station.
  \item \emph{Recall:} say the Bengali for to try a case, then the Bengali for a verdict, then say \emph{nālish} again
\end{itemize}

\begin{tcolorbox}[breakable,colback=teal!4,colframe=teal!35!black,title={Before you move on}]
What does \bn{নালিশ} mean? (\textquotedblleft{}A formal complaint\textquotedblright{}.) Say it once more.
\end{tcolorbox}
\section[\texorpdfstring{\bn{মীমাংসা} \bn{করা} (mīmāngsā kôrā)}{মীমাংসা করা (mīmāngsā kôrā)}]{\bn{মীমাংসা} \bn{করা} (mīmāngsā kôrā) — to settle a dispute; to resolve}

\label{lesson:BN-C342-mimangsa-kora}



\begin{quote}
Before the new one: say the Bengali for a verdict, then the Bengali for a formal complaint.
\end{quote}

\subsection*{You'll want to know: \bn{মীমাংসা} \bn{করা}}
\textbf{\bn{মীমাংসা} \bn{করা}} — \emph{mīmāngsā kôrā} — \textquotedblleft{}to settle a dispute; to resolve\textquotedblright{}.

\begin{quote}
\bn{দুই} \bn{পরিবার} \bn{ঝগড়াটার} \bn{মীমাংসা} \bn{করল।}
\end{quote}

\begin{quote}
\emph{dui pôribār jhôgṛāṭār mīmāngsā kôrlô.}
\end{quote}

\begin{quote}
The two families settled the quarrel.
\end{quote}

Past (he/she …): \bn{তিনি} \bn{মামলাটার} \bn{মীমাংসা} \bn{করলেন}, \emph{tini māmlāṭār mīmāngsā kôrlen}.

\begin{cousinweb}
\textbf{\bn{মীমাংসা} \bn{করা}} (\emph{mīmāngsā kôrā}) is Sanskrit \emph{mīmāṃsā}, deep reflection or investigation (Monier-Williams), with \bn{করা}. The Bengali sense, settling a question or a dispute, is not one the dictionary gives; the settled matter takes -\bn{এর}: \bn{ঝগড়ার} \bn{মীমাংসা}.
\end{cousinweb}

\begin{grammarlens}[title={What you've built}]
One more word for giving reasons and opinions.
\end{grammarlens}

\subsection*{Your turn}
\begin{itemize}
\raggedright
  \item \emph{Say it:} \emph{mīmāngsā kôrā}
  \item \emph{Say it:} \emph{mīmāngsā kôrā}, once more
  \item \emph{Say it:} \emph{dui pôribār jhôgṛāṭār mīmāngsā kôrlô.} — The two families settled the quarrel.
  \item \emph{Recall:} say the Bengali for a verdict, then the Bengali for a formal complaint, then say \emph{mīmāngsā kôrā} again
\end{itemize}

\begin{tcolorbox}[breakable,colback=teal!4,colframe=teal!35!black,title={Before you move on}]
What does \bn{মীমাংসা} \bn{করা} mean? (\textquotedblleft{}To settle a dispute; to resolve\textquotedblright{}.) Say it once more.
\end{tcolorbox}
